% =============================================================================
%  Full-ascent extension of the indirect-PMP (TPBVP) guidance — drafting source.
%
%  This snippet holds ONLY the genuinely-new material: the full-ascent extension
%  (drag-aware adjoints, angle-of-attack constraint + atmosphere exit, multi-arc
%  ascent with staging continuity, and the passive mass costate). The core
%  indirect-PMP formulation (objective, Hamiltonian, costate ODEs, control law,
%  transversality, penalty objective) already lives in the thesis, in
%  Thesis_Background_claude.tex, Section "Indirect Optimization of the Ascent
%  Trajectory" (\label{ssec:indirect_ascent}) — so it is NOT repeated here.
%
%  These four \subsubsection blocks are already INLINED into that section
%  (after "Penalty-Based Reformulation", before \section{Numerical Methods}).
%  This file is kept as the standalone drafting source; it is not \input by the
%  thesis. New \cite keys (gath2001optimization, pontani2013ascent,
%  pontani2019ascent) are defined in tpbvp_refs.bib for merging into the master
%  bibliography; the remaining keys reuse existing master-bib entries.
% =============================================================================

\subsubsection{Extension to the Atmospheric Arc: Drag-Aware Adjoints}
\label{sssec:fa_drag}

The formulation above applies the costate-based control only on the exo-atmospheric arcs, where aerodynamic forces are negligible. The same indirect framework can be carried through the atmospheric flight of the first stage by retaining aerodynamic drag in both the dynamics and the adjoint equations. Writing the drag specific force (deceleration) as
\begin{equation}
a_D = \frac{D}{m} = \frac{1}{2}\,\rho(h)\,v^{2}\,\frac{C_D\,S_{\mathrm{ref}}}{m} \ ,
\label{eq:fa_drag}
\end{equation}
with the exponential density of Eq.~\eqref{eq:exp_atmosphere}, drag enters the velocity dynamics as an additional $-a_D$ term in $\dot{v}$. Because $a_D$ depends on $r$ and $v$ but not on $\gamma$, and because the exponential atmosphere yields the closed-form partial derivatives $\partial a_D/\partial r = -a_D/H_\rho$ and $\partial a_D/\partial v = 2a_D/v$ (with $H_\rho$ the density scale height, written $H$ in Eq.~\eqref{eq:exp_atmosphere} and renamed here to avoid confusion with the Hamiltonian), the drag contribution to the adjoint equations~\eqref{eq:indirect_costates} is analytic and affects only $\dot{\lambda}_r$ and $\dot{\lambda}_v$. The extra terms, added to the right-hand sides of Eqs.~\eqref{eq:costate_r} and~\eqref{eq:costate_v}, are~\cite{gath2001optimization}:
\begin{equation}
\dot{\lambda}_r \mathrel{+}= -\,\lambda_v\,\frac{a_D}{H_\rho} \ , \qquad
\dot{\lambda}_v \mathrel{+}= +\,2\,\lambda_v\,\frac{a_D}{v} \ .
\label{eq:fa_drag_adjoint}
\end{equation}
A matching $a_D$ term is likewise retained in the Hamiltonian~\eqref{eq:indirect_hamiltonian}, so the transversality residual stays consistent with the drag-aware dynamics. In the vacuum limit $\rho\to 0$ these terms vanish and the drag-free equations of the preceding subsections are recovered exactly.

\subsubsection{Angle-of-Attack Constraint and Atmosphere Exit}
\label{sssec:fa_alpha}

Flying the powered first stage through the dense lower atmosphere requires limiting the angle of attack, both to respect aerodynamic load limits near maximum dynamic pressure and to keep the vehicle close to a gravity turn~\cite{Cremaschi2013TrajectoryOptimization}. The interior optimal control~\eqref{eq:indirect_control} is therefore clamped to an admissible band,
\begin{equation}
\alpha = \max\!\bigl(-\alpha_{\max},\,\min(\alpha_{\max},\,\alpha^{\star})\bigr) \ ,
\label{eq:fa_alpha_clamp}
\end{equation}
where $\alpha^{\star}$ denotes the unconstrained value given by Eq.~\eqref{eq:indirect_control} and $\alpha_{\max}$ is a prescribed bound. The constraint is enforced only while the vehicle remains within the atmosphere and is released once an atmosphere-exit criterion is met---here the dynamic pressure falling below a fixed threshold, $q < q_{\mathrm{exit}}$---so that the near-vacuum arc recovers the exact interior control~\eqref{eq:indirect_control}. Using the same exit criterion that governs fairing jettison and guidance activation keeps the notion of ``atmospheric flight'' consistent across the simulator~\cite{gath2001optimization}.

\subsubsection{Multi-Arc Ascent and Staging Continuity}
\label{sssec:fa_arcs}

With drag and the angle-of-attack constraint in place, the indirect law can steer the entire powered ascent rather than the upper stage alone. The trajectory is represented as an ordered sequence of arcs sharing a single continuous integration of the augmented state: a vertical rise at $\alpha\equiv 0$ up to the pitch-over instant, the first-stage powered arc terminated by propellant depletion (MECO), an instantaneous inert-mass separation (staging), an inter-stage ballistic coast, and finally the upper-stage thrust--coast--thrust arcs of the baseline formulation. The costates are seeded at the pitch-over point rather than at lift-off: the vertical rise carries no steering freedom ($\alpha\equiv 0$) and the homogeneous costate equations keep $\boldsymbol{\lambda}$ at zero through it, so fixing the normalised initial costates at the kick is both sufficient and numerically cleaner. Staging is a prescribed exogenous mass drop; because the reduced costate set carries no mass adjoint, the position and velocity costates and the Hamiltonian remain continuous across the staging and coast corners by the Weierstrass--Erdmann conditions, and no interior jump conditions are required~\cite{pontani2014particle,pontani2013ascent}.

\subsubsection{Mass Costate and Transversality}
\label{sssec:fa_masscostate}

Because the high mass-flow first stage consumes a large fraction of the lift-off mass, the mass adjoint $\lambda_m$ is retained on the full-ascent arcs for rigour. As the mass rate $\dot{m} = -T/(I_{sp}g_0)$ is independent of $r$, $v$ and $\gamma$, the partial $\partial H/\partial m$ produces no feedback into the steering costates, so $\lambda_m$ is a passive quadrature,
\begin{equation}
\dot{\lambda}_m = \lambda_v\!\left(\frac{T\cos\alpha}{m^{2}} - \frac{a_D}{m}\right)
+ \lambda_\gamma\,\frac{T\sin\alpha}{m^{2}v} \ ,
\label{eq:fa_masscostate}
\end{equation}
that neither alters the trajectory nor enlarges the set of unknowns. With no terminal mass cost the transversality condition is $\lambda_m(t_f)=0$, which, for a passive quadrature, is satisfied in closed form by subtracting the constant $\lambda_m(t_f)$ from the integrated history. Including $\lambda_m$ renders the Hamiltonian constant along each powered arc, confirming that the multi-arc solution is a consistent extremal of the full-ascent problem~\cite{pontani2013ascent,pontani2019ascent}.
